\title{Lecture 1: MDP Basics and Planning}
% \author{Qi Zhang (qzhang9@wpi.edu)}
% \date{Last Updated: August 2025}

\section{Markov Decision Processes: The Finite-Horizon Setting}
In reinforcement learning, the interaction between the agent and its environment is often formulated as a Markov Decision Process (MDP).
This note focuses on the {\em finite-horizon} setting where an MDP is specified by tuple $M:=(\mathcal{S},\mathcal{A}, P, R, H)$:
\begin{itemize}
    \item State space $\mathcal{S}$. This note only considers finite state spaces.
    \item Action space $\mathcal{A}$. This note only considers finite action spaces.
    \item Horizon $H$. This is a constant positive integer. Let $h\in[H]:=\{1,2,\ldots,H\}$ index the discrete timesteps.
    \item Transition function $P:\mathcal{S}\times\mathcal{A}\times[H]\to\Delta(\mathcal{S})$, where is the space of probability distributions over $\mathcal{S}$. $P_h(s'|s,a)$ is the probability of transiting to state $s'$ after taking action $a$ in state $s$.
    \item Reward function $R:\mathcal{S}\times\mathcal{A}\times[H]\to[0,1]$. $R_h(s,a)$ is the  immediate reward associated with taking action $a$ in state $s$ at timestep $h$. 
\end{itemize}

\subsection{Interaction protocol}
Starting in some state $s_1\in\mathcal{S}$, at each timestep $h\in[H]$, the agent takes an action $a_h \in \mathcal{A}$, obtains the immediate reward $r_h:=R_h(s_h, a_h)$, and observes the next state $s_{h+1}\in\mathcal{S}$ sampled from $P_h(s_h, a_h)$, or $s_{h+1} \sim P_h\left(s_h, a_h\right)$. The interaction record
\begin{align*}
    \tau:=\left(s_1, a_1, r_1, s_2, \ldots, s_H, a_H, r_H, s_{H+1}\right)
\end{align*}
is called an {\em episode}, which is a trajectory of length $H$.

In this note, we use notation $:=$ to denote assignment/definition, i.e., the left-hand side is defined to be equal to the right-hand side.

The process above can be iterated to form multiple episodes. We often consider the case where the initial state $s_1$ is fixed. More generally, $s_1$ is generated by sampling from a distribution $d_1\in \Delta(\mathcal{S})$.
When $d_1$ is of importance to the discussion, we include it as part of the MDP tuple, writing $M=(\mathcal{S},\mathcal{A}, P, R, H, d_1)$.


\section{Policy and Value}
A {\em policy} specifies how the actions are chosen.
A policy denoted as $\pi:\mathcal{S}\times[H]\to\mathcal{A}$ chooses actions deterministically based on the current state and the timestep, i.e., $a_h=\pi_h(s_h)$.
More generally, a policy denoted as $\pi:\mathcal{S}\times[H]\to\Delta(\mathcal{A})$ chooses actions stochastically denoted as $a_h\sim\pi_h(s_h)$, writing the probability as $\pi_h(a_h|s_h)$.
It is important to base the policy on the timestep $h$, because the optimal decision-making strategy may well depend on $h$ in this finite-horizon setting.
We therefore often write $\pi=\{\pi_h\}_{h=1}^H$.

Given MDP $M=(\mathcal{S},\mathcal{A}, P, R, H, d_1)$ and policy $\pi=\{\pi_h\}_{h=1}^H$ in place, one can explicitly write out the probability of generating any trajectory $\tau=\left(s_1, a_1, r_1, s_2, \ldots, s_H, a_H, r_H, s_{H+1}\right)$:
\begin{align*}
    \textstyle\Pr^{M,\pi}(\tau) = d_1(s_1)\pi_1(a_1|s_1)P_1(s_2|s_1,a_1)\cdots\pi_H(a_H|s_H)P_H(s_{H+1}|s_H,a_H)
\end{align*}
where rewards will be deterministically generated as $r_h=R_h(s_h,a_h)$.

The goal of the agent is to find a policy $\pi$ that maximizes the expected sum of rewards received in the future, or {\em values}.
In MDP $M$, following policy $\pi$ from state $s\in\mathcal{S}$ at timestep $h\in[H]$, the expected cumulative reward to the end is denoted as $V^{M,\pi}_h(s)$, i.e.,
\begin{align*}  V^{M,\pi}_h(s):=\E_{M,\pi}\left[\sum_{h^{\prime}=h}^H R_{h^{\prime}}\left(s_{h^{\prime}}, a_{h^{\prime}}\right) \mid s_h=s\right]
,
\end{align*}
which is called the state-value.
Similarly, the action-value (or Q-value) is defined as 
\begin{align*}
    Q^{M,\pi}_h(s, a):=\E_{M,\pi}\left[\sum_{h^{\prime}=h}^H R_{h^{\prime}}\left(s_{h^{\prime}}, a_{h^{\prime}}\right) \mid s_h=s, a_h=a\right]
\end{align*}
i.e., the expected cumulative reward by following policy $\pi$ in MDP $M$ from taking action $a$ in state $s\in\mathcal{S}$ at timestep $h\in[H]$.

When the MDP $M$ is clear from the context, we often drop superscript/subscript $M$ and write $V^\pi_h$, $Q^\pi_h$, $\Pr^{\pi}$, $\E_\pi$, etc.

The goal of the agent is to find a policy $\pi$ that maximizes $V_1^\pi\left(s_1\right)$ if the initial state $s_1$ is fixed or $\E_{s_1 \sim d_1} [V_1^\pi\left(s_1\right)]$ if the initial state $s_1$ is sampled from distribution $d_1$.

\section{Planning in MDPs}
{\em Planning} refers to the problem of computing a value-maximizing policy given the full MDP specification $M=(\mathcal{S},\mathcal{A}, P, R, H)$.
A related problem is {\em policy evaluation}, which aims to obtain the values of a given policy. 
\subsection{Bellman equation for policy evaluation}
Policy evaluation is the problem of finding the values ($V^\pi_h$ and/or $Q^\pi_h$) of a given policy $\pi$.
In the context of planning, we also assume the full knowledge of the MDP specification $M$.
By definition, the values can be computed by expanding the expectation, e.g., 
\begin{align*}
    V^\pi_1(s) = \sum_{\tau:s_1=s} \left[\textstyle\Pr^\pi(\tau|s_1=s)\cdot\sum_{h=1}^{H}r_h\right]
\end{align*}
where $\tau:s_1=s$ enumerates all trajectories with the first state $s_1=s$ and $\Pr^\pi(\tau|s_1=s)=\pi_1(a_1|s_1=s)P_1(s_2|s_1=s,a_1)\cdots\pi_H(a_H|s_H)P_H(s_{H+1}|s_H,a_H)$ is the probability of getting such a trajectory $\tau$ by following $\pi$ starting in $s_1=s$.
This approach is highly inefficient because it enumerates the trajectories. In particular, there are in total $|\mathcal{S}\times\mathcal{A}|^H$ trajectories of length $H$, which is exponentially large in horizon $H$.

A more efficient way to do policy evaluation is based on the principles of dynamic programming.
By definition, the values can be computed recursively via the following {\em Bellman equations}: letting $V^\pi_{H+1}(s)\equiv0$, $\forall s\in\mathcal{S},a\in\mathcal{A},h\in[H]$,
\begin{align}
\begin{split}\label{eq:Bellman}
& V_h^\pi(s)=\E_{a \sim \pi_h(s)} [Q_h^\pi(s, a)],\\
& Q_h^\pi(s, a)=R_h(s, a)+\E_{s' \sim P_h( s, a)}[V_{h+1}^\pi(s')]
\end{split}
\end{align}
where the recursion proceeds backward in time as $h=H,H-1,\ldots,1$, so the computation is linear in horizon $H$.

It is often convenient to rewrite Eq.~\eqref{eq:Bellman} in a matrix-vector form.
Since $\mathcal{S}$ and $\mathcal{A}$ are assumed to be finite, upon fixing an arbitrary order of states and actions, we can treat the relevant quantities as matrices or vectors of proper shapes.
Specifically, Eq.~\eqref{eq:Bellman} can be rewritten as
\begin{align*}
    V_h^\pi =&\left(\pi_h \odot Q_h^\pi\right)\mathbf{1}
    \quad \text{with }V_h^\pi\in\R^{|\mathcal{S}|}, \pi_h ,Q_h^\pi \in \R^{|\mathcal{S}| \times |\mathcal{A}|}, \text{ all-one vector } \mathbf{1} \in\R^{|\mathcal{A}|}, 
    \\
    Q_h^\pi =& R_h + P_h V^\pi_h \quad \text{with }
    Q_h^\pi, R_h \in \R^{|\mathcal{S} \times \mathcal{A}|},
    P_h\in\R^{|\mathcal{S} \times \mathcal{A}|\times|\mathcal{S}|}, V_{h+1}^\pi\in\R^{|\mathcal{S}|}
    .
\end{align*}
Note that we treat $Q_h^\pi$ above as a matrix in the first equation and a vector in the second equation.

\subsection{Bellman optimality equation}
With a slight modification of replacing the $\E_{a \sim \pi_h(s)}$ in Bellman equations with the greedy action $\max_{a\in\mathcal{A}}$, we obtain the {\em Bellman optimality equations}:
letting $V^*_{H+1}(s)\equiv0$, $\forall s\in\mathcal{S},a\in\mathcal{A},h\in[H]$,
\begin{align}
\begin{split}\label{eq:Bellman_optimality}
& V_h^*(s):=\max_{a\in\mathcal{A}}  Q_h^*(s, a),\\
& Q_h^*(s, a):=R_h(s, a)+\E_{s' \sim P_h( s, a)}[V_{h+1}^*(s')]
\end{split}
\end{align}
which recursively define quantities $V^*_h$ and $Q^*_h$.
Consider the policy that is acting greedy with respect to $Q^*_h$, i.e.,
\begin{align}\label{eq:optimal_policy}
    \pi^*_h(s) := \argmax_{a\in\mathcal{A}} Q^*_h(s,a) \quad\forall s\in\mathcal{S}
\end{align}
We have Proposition \ref{proposition:optimality} showing that $\pi^*$ is an optimal policy with its values equal to $V^*$ and $Q^*$.
\begin{proposition}
\label{proposition:optimality}
For $V^*$ and $Q^*$ defined in \eqref{eq:Bellman_optimality} and policy $\pi^*$ defined in \eqref{eq:optimal_policy}, we have
\begin{align*}
    \quad V_h^*(s)=V^{\pi^*}_h(s)=\max _\pi V_h^\pi(s), \quad Q_h^*(s, a)=Q^{\pi^*}_h(s,a)=\max _\pi Q_h^\pi(s, a)
    \qquad\forall(s, a, h)
\end{align*}
where the maximization is taken over all policies $\pi=\{\pi_h\}_{h=1}^H$.
\end{proposition}
\begin{proof}
We prove it by induction over $h=H+1,H,\ldots,1$.

\paragraph{Base case.} For $h=H+1$, we have $V^*_{H+1}(s)\equiv0=V^{\pi^*}_{H+1}(s)=\max _\pi V_{H+1}^\pi(s)$ because $V^\pi_{H+1}(s)\equiv0$ for any $\pi$.

~\\
Assume the statement in the proposition holds for $h+1$, i.e.,
$\forall s\in\mathcal{S}, a\in\mathcal{A}$
\begin{align*}
    \quad V_{h+1}^*(s)=V^{\pi^*_{h+1:H}}_{h+1}(s)=\max _{\pi_{h+1:H}} V_{h+1}^{\pi_{h+1:H}}(s), \quad Q_{h+1}^*(s, a)=Q^{\pi^*_{h+1:H}}_{h+1}(s,a)=\max_{\pi_{h+1:H}} Q_{h+1}^\pi(s, a)
    .
\end{align*}
Note here we explicitly write $V^\pi_{h+1}\equiv V^{\pi_{h+1:H}}_{h+1}$ (and for $Q$ as well) to emphasize the fact that the value function for timestep $h+1$ does not depend on policy before that timestep.
The induction proceeds to $h\in[H]$ in the following three steps.
\paragraph{1) $Q^*_h$ is optimal.}
We have, $\forall (s,a)$, 
\begin{align*}
&\max_\pi ~ Q_h^\pi(s, a) \\
=&\max _\pi ~ R_h(s, a)+\E_{s'\sim P_h(s, a)} [V_{h+1}^\pi(s')] \qquad&&\text{(1a)}\\
=&R_h(s, a)+\max _\pi \E_{s'\sim P_h(s, a)} [V_{h+1}^\pi(s')] \\
=&R_h(s, a)+\E_{s'\sim P_h(s, a)} \left[V_{h+1}^{\pi^*_{h+1:H}}(s')\right] \qquad&&\text{(1b)}\\
=&R_h(s, a)+\E_{s'\sim P_h(s, a)} [V_{h+1}^*(s')] \qquad&&\text{(Induction hypothesis)}\\
=&Q_h^*(s, a) \qquad&&\text{(1c)}
\end{align*}

\paragraph{2) $V^*_h$ is optimal.}
We have, $\forall s$, 
\begin{align*}
\max_\pi ~V_h^\pi(s)
=&\max_\pi ~\sum_a \pi_h(a | s) Q_h^\pi(s, a) \\
=&\max_{\pi_h, \pi_{h+1: H}} ~\sum_a \pi_h(a | s) Q_h^{\pi_{h+1: H}}(s, a)\\ 
=&\max_{\pi_h} ~\sum_a \pi_h(a | s) Q_h^*(s, a)  \qquad&&\text{(1d)}  \\
=&\max_a ~Q^*_h(s,a) \\
=&V^*_h(s) \qquad&&\text{(1e)}
\end{align*}

\paragraph{3) $\pi^*_{h:H}$ is optimal.}
We have, $\forall s$,
\begin{align*}
V_h^{\pi_{h:H}^*}(s)
=& Q_h^{\pi^*_{h+1: H}}(s, \bar{a}) \qquad &&\text {(where $ \bar{a}:=\pi_h^*(s)=\argmax_{a\in\mathcal{A}} Q_h^*(s, a)$)} \\
=& R_h(s, \bar{a})+\mathbb{E}_{s' \sim P_h(s, \bar{a})}\left[V_{h+1}^{\pi_{h+1: H}^*}(s')\right]\\
=& R_h(s, \bar{a})+\mathbb{E}_{s' \sim P_h(s, \bar{a})}\left[V_{h+1}^{*}(s')\right] \qquad&&\text{(1f)}\\
= &Q_h^*(s, \bar{a}) \\
= &\max_a Q_h^*(s, a) = V^*_h(s)
\end{align*}
The proof completes because steps (1,2,3) finish the induction on $h$.
\end{proof}
